% Figura: experimento de Cavendish y argumento de la cáscara esférica.
\begin{figure}[ht]
\centering
\tikzset{
  parche/.style={draw=azulprof, line width=0.6pt, fill=azulclaro!70, fill opacity=0.9},
  parche2/.style={draw=acento, line width=0.6pt, fill=acento!25, fill opacity=0.9},
  vn/.style={-{Stealth[length=5pt,width=4.5pt]}, line width=1.1pt}}
\begin{minipage}[t]{0.50\textwidth}\centering
\begin{tikzpicture}[baseline=(current bounding box.north), x=0.82cm, y=0.82cm]
  % (a) cáscara de dos hemisferios con esfera interior unida por un hilo conductor
  \def\Re{1.9}\def\Ri{0.6}
  \shade[ball color=azulclaro!55, opacity=0.35] (0,0) circle (\Re);
  \draw[azulprof, line width=0.8pt] (0,0) circle (\Re);
  \draw[azulprof, line width=0.8pt] (-\Re,0) arc (180:360:{\Re} and 0.45);
  \draw[azulprof, line width=0.5pt, densely dashed] (-\Re,0) arc (180:0:{\Re} and 0.45);
  \shade[ball color=azulacero!70, opacity=0.95] (0,0) circle (\Ri);
  \draw[azulprof, line width=0.7pt] (0,0) circle (\Ri);
  % hilo conductor vertical con collar aislante en la cáscara y terminal exterior
  \draw[grispizarra, line width=1.0pt] (0,\Ri) -- (0,2.55);
  \draw[azulprof, line width=0.6pt, fill=white] (0,\Re) ellipse (0.13 and 0.05);
  \fill[acento] (0,2.55) circle (2.0pt);
  % rótulos con líneas de llamada rectas
  \node[font=\footnotesize, anchor=west] at (0.35,2.55) {hilo conductor};
  \draw[grispizarra, line width=0.3pt] (-2.65,0.45) -- (-1.85,0.45); \node[font=\footnotesize, anchor=east] at (-2.7,0.45) {hemisferios};
  \draw[grispizarra, line width=0.3pt] (2.7,-0.9) -- (1.5,-0.9); \node[font=\footnotesize, anchor=west] at (2.75,-0.9) {cáscara};
  \draw[grispizarra, line width=0.3pt] (2.7,0.0) -- (0.62,0.0); \node[font=\footnotesize, anchor=west] at (2.75,0.0) {esfera interior};
  \draw[verdeteo, <->, line width=0.5pt] (0,0) -- (-35:\Ri) node[pos=0.7, below=1pt, font=\scriptsize, text=verdeteo] {$a$};
\end{tikzpicture}\\[2pt]{\small\scshape (a)}
\end{minipage}\hfill
\begin{minipage}[t]{0.48\textwidth}\centering
\begin{tikzpicture}[baseline=(current bounding box.north), x=1.05cm, y=1.05cm]
  % (b) cono desde un punto interior P que recorta dos elementos opuestos de la cáscara
  \def\Rr{2.0}
  \draw[azulprof, line width=0.9pt, name path=circ, fill=azulfondo, fill opacity=0.5] (0,0) circle (\Rr);
  \coordinate (O) at (0,0); \fill[azulprof] (O) circle (1.2pt); \node[font=\scriptsize, anchor=north east] at (O) {$C$};
  \coordinate (P) at (0.7,-0.25);
  \def\dd{24}\def\ap{9}
  \foreach \dir/\nm in {\dd-\ap/a,\dd+\ap/b,\dd/m}{
    \path[name path=ray\nm] (P) -- ++({\dir}:5);
    \path[name path=rev\nm] (P) -- ++({\dir+180}:5);
    \path[name intersections={of=circ and ray\nm, by=A\nm}];
    \path[name intersections={of=circ and rev\nm, by=B\nm}];}
  \fill[azulclaro, opacity=0.40] (P) -- (Aa) -- (Ab) -- cycle;
  \fill[azulclaro, opacity=0.40] (P) -- (Ba) -- (Bb) -- cycle;
  \draw[azulacero, line width=0.45pt] (Aa) -- (Ba) (Ab) -- (Bb);
  \draw[azulprof, line width=0.8pt] (Bm) -- (Am);
  \draw[acento, line width=2.8pt] (Aa) -- (Ab);
  \draw[acento, line width=2.8pt] (Ba) -- (Bb);
  \draw[verdeteo, line width=0.6pt, densely dashed] (0,0) -- (Am) (0,0) -- (Bm);
  \fill[azulprof] (P) circle (1.5pt) node[below=2pt, font=\small] {$P$};
  \node[font=\scriptsize, text=azulprof, anchor=south east] at ($(P)!0.55!(Am)$) {$r_2$};
  \node[font=\scriptsize, text=azulprof, anchor=north west] at ($(P)!0.55!(Bm)$) {$r_1$};
  \node[font=\small, text=acento, anchor=south west] at ($(Am)+(0.1,0.12)$) {$\dif S_2$};
  \node[font=\small, text=acento, anchor=north east] at ($(Bm)+(-0.12,-0.2)$) {$\dif S_1$};
  % ángulo alpha entre la cuerda y el radio en cada extremo (arcos dibujados a mano)
  \draw[verdeteo, line width=0.55pt] let \p1=(Am), \n1={atan2(\y1,\x1)+180} in
     (Am) ++(\n1:0.45) arc[start angle=\n1, end angle=\dd+180, radius=0.45];
  \draw let \p1=(Am), \n1={atan2(\y1,\x1)+180} in
     node[font=\scriptsize, text=verdeteo] at ($(Am)+({(\n1+\dd+180)/2}:0.78)$) {$\alpha$};
  \draw[verdeteo, line width=0.55pt] let \p1=(Bm), \n1={atan2(\y1,\x1)+180} in
     (Bm) ++(\n1:0.45) arc[start angle=\n1, end angle=\dd, radius=0.45];
  \draw let \p1=(Bm), \n1={atan2(\y1,\x1)+180} in
     node[font=\scriptsize, text=verdeteo] at ($(Bm)+({(\n1+\dd)/2}:0.78)$) {$\alpha$};
\end{tikzpicture}\\[2pt]{\small\scshape (b)}
\end{minipage}
\caption{(a)~Experimento de Cavendish: una esfera conductora interior se conecta, mediante un hilo, a una cáscara formada por dos hemisferios, y se cargan; al retirar el hilo y separar los hemisferios,
la esfera interior resulta descargada. (b)~Argumento geométrico: un cono de ángulo sólido $\dOm$ con vértice en un punto interior $P$ recorta dos elementos, $\dif S_1$ y $\dif S_2$,
a distancias $r_1$ y $r_2$ de $P$; ambos forman con la cuerda el mismo ángulo $\alpha$ (el triángulo formado por los dos radios y la cuerda es isósceles).}
\label{fig:cavendish}
\end{figure}
